\chapter{Introduction}
\label{chapter:intro}

\pagenumbering{arabic}

In recent years, deep learning and computer vision methods have made significant advancements in the analysis 
of dynamic scenes and their applications in autonomous systems. The development of fields such as autonomous 
vehicles and robotics depends on solving tasks like object detection, object tracking, segmentation, and 
trajectory prediction. Specifically, \ac{avs} revolutionize transportantion by saving drivers' time, improving safety and
possibly mitigating congestion on the roads. Current autonomous driving systems consist of multiple sequential tasks such as perception,
tracking, prediction, planning and control \cite{chen2024end}. Perception task is responsible for detecting and understanding
the objects within the environment including their categorization, location and size estimation. The tracking objective
uses the detections from perception to find the relationships between objects across the frames in scene. The prediction 
modules then process this tracking information in order to forecast future positions. The predicted positions enter the
planning stage that generates the most feasible and optimal path in terms of safety, energy consuption and following traffic rules.
Last but not least, controlling task then prepares actuator commands which execute the planned motion.

\section{Motivation}
Real traffic is highly dynamic and unpredictable: other drivers may behave erratically, and sensors provide 
noisy or incomplete data. Diverse maneuvers, complex interactions, and inherent errors in sensory information 
make prediction challenging, yet solving these challenges is essential for having trustworthy \ac{avs}. 
In addition, autonomous systems operate under strict constraints on latency and reliability \cite{huang2022survey}. 
Predictions must be computed in real time on embedded hardware and remain accurate over long horizons to be useful 
for driving in real-life scenarios. Industry recognizes the need for fast execution, uncertainty handling, and generalization 
in prediction models to ensure robust operation in varied conditions \cite{Ettinger2021WOMD,Wayformer2022,ISO21448}. 
Reliable trajectory prediction is therefore critical for safe autonomous driving and motivates intensive 
research into methods that can forecast the multi-agent future with practical efficiency. For these reasons, this 
thesis focuses on trajectory prediction, exploring how current state-of-the-art models could be modified or extended
to achieve better accuracy and efficiency.


\section{Problem Formulation}
Trajectory prediction is the task of forecasting the future motion of the traffic participants that surround
the ego vehicle. Its input is not a trajectory in the purely geometric sense, but rather the past states of
these participants as they are captured in the environmental information collected by sensors,
such as camera images, LiDAR point clouds, radar returns or HD maps providing lane geometry and traffic-rule
priors. However, in this thesis only camera and LiDAR is employed. A model that relies on a single source 
is unimodal, whereas a model that exploits several of them at
once is considered multi-modal. In the latter case the measurements are heterogeneous in dimensionality,
sampling rate and coordinate frame, so the features extracted from the individual sources have to be fused
into a unified representation of the scene before any forecasting takes place \cite{madjid2025survey}. The
output is the forecasted trajectory, that is, a sequence of future states of the observed agents over a fixed
prediction horizon.

Formally, trajectory prediction can be stated as the problem of constructing a mapping $\Phi$ from the
observed evidence $\mathcal{I}$ about the scene to the future states $\mathcal{O}$ of the traffic
participants, from the current time up to a prediction horizon $T$ \cite{hu2025survey}:
\begin{equation}
    \begin{aligned}
        \Phi &: \mathcal{I} \rightarrow \mathcal{O}, \\
        \mathcal{O} &= \left\{ S^{0}, S^{1}, \dots, S^{T} \right\},
    \end{aligned}
    \label{eq:traj-pred}
\end{equation}
where $S^{0}$ is the current state and $S^{1}$ to $S^{T}$ are the states at the individual future time
steps, each of them collecting the states of all $N$ observed participants. Depending on the available
modalities and on the downstream task, a single state can take different forms, for example
\begin{equation}
    S^{t} = \left\{ x^{t}, y^{t} \right\} \quad \text{or} \quad
    S^{t} = \left\{ x^{t}, y^{t}, l^{t}, w^{t}, \theta^{t} \right\},
    \label{eq:state}
\end{equation}
where $x^{t}$ and $y^{t}$ are the coordinates of the centre point of the object, $l^{t}$ and $w^{t}$ its
length and width, and $\theta^{t}$ its orientation \cite{hu2025survey}. In trajectory prediction the state is
usually reduced to the planar position, so that $\mathcal{O}$ corresponds to the set of forecasted
trajectories of the participants in the scene. Since the future of a traffic scene is inherently uncertain,
the mapping in \eqref{eq:traj-pred} is in practice often extended to output several plausible trajectories
together with their confidences, allowing the planning stage to reason over the alternatives instead of
committing to a single hypothesis \cite{huang2022survey}. The following subsections summarise the main
challenges that this formulation raises.

\subsection{Data Fusion}
The information describing the scene comes from sources that are complementary but incompatible. Cameras
provide dense semantic detail without depth and degrade at night, in glare or in poor lighting
conditions, LiDAR provides accurate geometry but suffers in rain etc.
These sources can be combined at the level of raw measurements (early), of learned features (intermediate) or of the individual predictions (late).
Since every added modality costs latency and introduces noise the network is prone to overloading with parameters or can overspecialize 
on the most informative source and lose accuracy once the weather or illumination degrades it \cite{madjid2025survey}.

\subsection{Dynamic Environments}
Traffic is nonstationary, as density, weather, illumination and road layout change, so that the conditions
met during deployment need not resemble those seen during development. The future is moreover not unique,
since the same observed history is usually compatible with several maneuvers, and this aleatoric uncertainty
originated in the erratic decisions of other drivers does not diminish with additional
data. Therefore the task requires a set of socially compliant
trajectories produced in time for the planning stage \cite{madjid2025survey}.

\subsection{Interaction Modeling}
The motions of traffic participants are not independent and traffic signals constrain all of them, so that the joint
future does not split into separate per-agent distributions. Capturing the agent to agent and the agent
to map relations requires representations able to operate on a variable and unordered number of entities,
which is why structured formulations such as graphs and attention mechanisms prevail. Where these relations
are modelled poorly, forecasts that appear plausible for each agent in isolation become inconsistent at the
level of the scene comprising of multiple entities \cite{huang2022survey}.

\subsection{Real-World Constraints}
Prediction is one stage of a control loop closing at a fixed rate, which leaves it a budget of tens of
milliseconds on embedded hardware limited in memory and power. Its cost grows with the number of
participants and becomes quadratic once their interactions are modelled, so the least favourable case is
dense traffic, precisely where prediction matters most. Model capacity therefore has to be weighed against
computational efficiency, as accuracy alone does not make a model deployable
\cite{Wayformer2022,huang2022survey}.


\section{Thesis Outline}



